\documentclass[conference,spanish]{IEEEtran}
\usepackage[utf8]{inputenc}
\usepackage[T1]{fontenc}
\usepackage{babel}
\usepackage{hyperref}
\usepackage{booktabs}
\usepackage{tabularx}
\usepackage{array}
\usepackage{ragged2e}
\usepackage{amsmath}

\newcolumntype{Y}{>{\RaggedRight\arraybackslash}X}
\newcolumntype{P}[1]{>{\RaggedRight\arraybackslash}p{#1}}

\title{Propuesta de proyecto: SAFESITE V2\\
\small Hito 1 -- Problema, ruta, datos y medidas de desempeño}

\author{
\IEEEauthorblockN{William David Lozano Julio \hspace{0.4em}--\hspace{0.4em} T00078475 \\
Cristian Matteo Benitez Montes \hspace{0.4em}--\hspace{0.4em} T00078658}
\IEEEauthorblockA{Curso de Inteligencia Artificial -- Universidad de Cartagena\\
Profesor: Edwin Puerta}
}
\date{3 de noviembre de 2026}

\begin{document}
\maketitle

\section{Problema y usuario}
En las empresas de construcción, el supervisor de seguridad y salud en el trabajo (HSE) necesita saber quién ingresa a la obra, si está registrado en la nómina, cuándo inicia y termina su jornada y si utiliza el Equipo de Protección Personal (EPP) requerido para su puesto. Actualmente, estas tareas pueden realizarse mediante registros manuales y controles separados. Esto dificulta consolidar asistencia, puntualidad, horas trabajadas e incidentes de seguridad en un mismo flujo.

Se propone \textbf{SAFESITE V2}, una evolución del agente original de detección de EPP hacia un gestor y verificador de nómina LIVE. El usuario principal es el supervisor HSE o encargado de obra. El sistema debe identificar al trabajador mediante badge/QR o selección manual, validar su registro, marcar entrada y salida, calcular horas trabajadas y generar reportes agregados en PDF y WhatsApp. No se utilizará reconocimiento facial ni se tomarán decisiones salariales irreversibles sin revisión humana.

\section{Ruta de IA elegida y arquitectura}
La ruta principal es la \textbf{E: agente con herramientas}. El Agente 1, de Nómina e Identificación, utiliza un LLM con herramientas para consultar empleados y turnos, registrar asistencia, verificar puntualidad, calcular horas y solicitar reportes. Los cálculos numéricos y las reglas laborales permanecen deterministas para facilitar la auditoría.

La Ruta F, visión por computador, funciona como apoyo mediante el Agente 2 de Verificación de EPP. Este agente utiliza YOLOv8s y confirmación temporal de $K=3$ frames para comprobar el EPP requerido por el puesto. El orden es secuencial:
\[
\text{identificación y nómina} \rightarrow \text{entrada} \rightarrow \text{EPP} \rightarrow \text{salida} \rightarrow \text{reporte}.
\]

Si el trabajador no existe en la nómina, se registra el intento no autorizado y se genera una alerta. Si existe, el Agente 1 registra la entrada y el Agente 2 verifica el EPP. La falta de EPP genera un incidente y una observación, pero no sustituye la validación de identidad ni bloquea automáticamente la entrada. El proyecto se organizará con CRISP-ML(Q): comprensión del negocio y los datos, preparación, modelado, evaluación, despliegue y control continuo de calidad.

\section{Datos, origen y licencia}
La solución utilizará dos fuentes de datos. Para el Agente 2 se empleará el dataset oficial \emph{Construction-PPE} de Ultralytics, con 1,416 imágenes y 11 clases relacionadas con EPP presente y ausente. La documentación oficial indica licencia AGPL-3.0; se conservará la atribución y se respetarán sus condiciones de uso \cite{ultralytics_dataset}. La copia disponible contiene 1,413 imágenes por la expiración de tres URL del CDN. Se utilizarán los splits preparados en \texttt{dataset\_v2/}: 1,130 imágenes de entrenamiento, 143 de validación y 140 de prueba local.

Para el Agente 1 se generarán datos sintéticos de nómina con semilla 42. Incluirán empresas, puestos, turnos, empleados pseudonimizados, asistencias, incidentes, períodos de nómina y reglas configurables. No se usarán nombres reales, rostros, cédulas ni salarios reales. La identificación se simulará mediante códigos badge/QR y casos controlados de credencial válida, inválida o duplicada.

\section{Línea base y solución propuesta}
La línea base principal será un LLM con prompt único y sin herramientas, evaluado en las mismas tareas de nómina. La solución será el Agente 1 con herramientas deterministas y estado de asistencia, conectado secuencialmente al Agente 2 YOLOv8s. Esta comparación permitirá medir si el uso de herramientas mejora la ejecución de tareas frente a una respuesta conversacional sin acceso a datos ni acciones externas. La visión EPP se comparará adicionalmente con el detector base sin confirmación temporal.

\begin{table}[htbp]
\centering
\caption{Medidas y umbrales definidos antes de experimentar}
\label{tab:metricas}
\scriptsize
\begin{tabularx}{\columnwidth}{P{0.29\columnwidth}Y P{0.19\columnwidth}}
\toprule
\textbf{Medida} & \textbf{Fórmula o procedimiento} & \textbf{Aceptación} \\
\midrule
Éxito de tarea del Agente 1 & Tareas de nómina completadas correctamente / tareas intentadas; mínimo 20 casos & $\geq 0.80$ \\
Estado de asistencia & F1 para presente, tardanza, bloqueado y ausente & $\geq 0.85$ \\
Errores de herramientas & Llamadas incorrectas o fallidas / llamadas totales & $\leq 0.05$ \\
Error de horas & Error absoluto entre horas calculadas y valor esperado & $\leq 1$ min \\
Recall de incidente EPP & Escenas con falta alertadas / escenas con falta & $\geq 0.70$ \\
Reporte y respuesta & Reportes PDF correctos; tiempo por tarea completa & $\geq 0.95$; $<5$ s \\
\bottomrule
\end{tabularx}
\end{table}

Cada medida se probará con escenarios normales, límites y fallas esperadas: badge inexistente, empleado duplicado, entrada sin salida, cambio de turno, frame inválido, falta de EPP y cierre de jornada con varias personas. Cuando exista aleatoriedad se realizarán al menos cinco repeticiones, con media y desviación estándar. Los logs registrarán percepción, herramienta invocada, decisión, acción y timestamp.

\section{Resultado esperado}
SAFESITE V2 debe producir un flujo reproducible de control de acceso, asistencia y cumplimiento de EPP, con separación clara entre el razonamiento con herramientas del Agente 1 y la percepción visual del Agente 2. La aprobación de esta propuesta depende de que el alcance, los datos, la línea base y los umbrales sean viables antes de iniciar las pruebas formales.

\begin{thebibliography}{9}
\bibitem{ultralytics_dataset}
Ultralytics, ``Construction-PPE Dataset,'' documentación oficial, 2026. Disponible en: \url{https://docs.ultralytics.com/datasets/detect/construction-ppe}
\end{thebibliography}

\end{document}
